\section*{Hoofdstuk \ref{ch:b1:clastic-sediments} --- Klastische sedimenten en hun milieus}

\begin{solution}{exo:b1:clastic-sediments:1}
$\varphi = -\log_2 d$: \qty{0.18}{mm} geeft $2.47$, een fijn zand (2 tot 3);
\qty{0.75}{mm} geeft $0.42$, een grof zand (0 tot 1); \qty{3}{mm} geeft
$-1.58$, een granule ($-1$ tot $-2$); \qty{0.012}{mm} geeft $6.38$, een fijne silt (6
tot 7); \qty{150}{mm} geeft $-7.23$, een kei ($-6$ tot $-8$).
\end{solution}

\begin{solution}{exo:b1:clastic-sediments:2}
$M_z = (1.05 + 1.71 + 2.43)/3 = 1.73\,\varphi$, dat is $2^{-1.73} =
\qty{0.30}{mm}$, een middelgrof zand. $\sigma_I = (2.43 - 1.05)/4 + (2.98 - 0.62)/6.6
= 0.345 + 0.358 = 0.70\,\varphi$: matig goed gesorteerd (0.50 tot 0.71), net
onder de grens van matig gesorteerd.
\end{solution}

\begin{solution}{exo:b1:clastic-sediments:3}
$q_b = (1 - 0.4) \times 0.02 \times (\num{0.01}/60)/2 = \qty{1.0e-6}{m^2.s^{-1}}$
aan vaste korrels. In een dag is dat $\num{1.0e-6} \times 86400 = \qty{0.0864}{m^3}$ per
meter kam, of $0.0864 \times 2650 = \qty{229}{kg}$ kwarts.
\end{solution}

\begin{solution}{exo:b1:clastic-sediments:4}
(a)~50 knopen zijn $50 \times 1852/3600 = \qty{25.7}{m.s^{-1}}$; 12 knopen,
\qty{6.17}{m.s^{-1}}. (b)~$\qty{100000}{m}/\qty{6.17}{m.s^{-1}} =
\qty{16200}{s}$, ongeveer \qty{4.5}{h}. (c)~480 mijl zijn \qty{772}{km}, 400 mijl
\qty{644}{km}.
\end{solution}

\begin{solution}{exo:b1:clastic-sediments:5}
(a)~$(\varphi_{84} - \varphi_{16})/2 = 2 \times 0.9945\,\sigma/2 =
0.9945\,\sigma$, en $(\varphi_{95} - \varphi_5)/3.3 = 2 \times
1.6449\,\sigma/3.3 = 0.9969\,\sigma$: beide schatten $\sigma$ binnen een derde
procent. (b)~Een echt sediment hoeft in phi-eenheden niet normaal te zijn: een kleine
bijmenging van veel grovere of veel fijnere korrels verandert $\varphi_5$ of
$\varphi_{95}$ sterk en $\varphi_{16}$ en $\varphi_{84}$ weinig, zodat de tweede schatting
groeit terwijl de eerste dat niet doet. De twee verschillen het meest voor mengsels, zoals
een zand met wat grind en modder, of een bimodale afzetting.
\end{solution}

\begin{solution}{exo:b1:clastic-sediments:6}
(a)~$\sin\theta = T/L = 1.5/15 = 0.1$, dus $\theta = \ang{5.74}$. (b)~$a =
c\tan\theta = 0.5 \times 0.1005 = \qty{0.050}{cm}$ per minuut, of
\qty{3.0}{cm} per uur. (c)~Een bodem die met centimeters per uur wordt opgehoogd, ontvangt
zand dat zeer snel uit suspensie valt: een zwaar beladen stroming die vertraagt, zoals een
afnemende vloed of de staart van een \omterm{def:b1:clastic-sediments:turbidite}{troebelingsstroom}.
\end{solution}

\begin{solution}{exo:b1:clastic-sediments:7}
(a)~$(350 + 20 + 5 + 340 + 15 + 30)/6 = 760/6 = 126.7$: een zuidoostelijke
richting, hoewel elke meting binnen \ang{20} van het noorden ligt. (b)~$S =
\sum\sin\alpha_i = 0.672$, $C = \sum\cos\alpha_i = 5.692$, dus $\bar\alpha =
\operatorname{atan2}(0.672, 5.692) = \ang{6.7}$, en $\bar R =
\sqrt{0.672^2 + 5.692^2}/6 = 5.732/6 = 0.955$. (c)~Het rekenkundig gemiddelde behandelt
350 en 005 als ver uit elkaar, terwijl zij op de cirkel \ang{15} uit elkaar liggen:
\omterm{def:b1:spherical-astronomy:horizontal}{azimuten} moeten als eenheidsvectoren worden gemiddeld. De grote $\bar R$ laat zien dat de
metingen strak rond \ang{7} zijn gegroepeerd.
\end{solution}

\begin{solution}{exo:b1:clastic-sediments:8}
A heeft meer dan \qty{5}{\percent} klei: onrijp, en compositioneel eveneens onrijp, met
zijn \omterm{def:b1:rock-forming-minerals:groups}{veldspaat} en gesteentefragmenten; snel afgezet nabij de bron, op een aanspoelwaaier
of een uiterwaarde, of door een \omterm{def:b1:clastic-sediments:turbidite}{troebelingsstroom}. B heeft minder dan \qty{5}{\percent}
klei en is goed gesorteerd ($\sigma_I < 0.5\,\varphi$), maar haar korrels zijn
subhoekig ($\rho < 3$): rijp, bijvoorbeeld een rivierbank of een plat zand dat door
stromingen is uitgezeefd. C heeft geen klei, is zeer goed gesorteerd en haar korrels zijn
afgerond ($\rho > 3$): superrijp, en met \qty{98}{\percent} kwarts ook compositioneel
rijp, zoals op een strand of in een woestijnduinveld.
\end{solution}

\begin{solution}{exo:b1:clastic-sediments:9}
(a)~$1.5/\sqrt{9.81 \times 0.3} = 0.87$; $1.0/\sqrt{9.81 \times 2} = 0.23$.
(b)~$h = U^2/g = 1.5^2/9.81 = \qty{0.23}{m}$; $U = \sqrt{9.81 \times 2} =
\qty{4.43}{m.s^{-1}}$. (c)~Het Froude-getal groeit naarmate de diepte bij een gegeven
snelheid afneemt: dunne, snelle waterlagen, zoals de terugstroom van een golf op een
strand, bereiken gemakkelijk Froude-getallen nabij 1, terwijl een rivier van enkele meters
diep enkele meters per seconde nodig zou hebben, wat weinig rivieren behalve tijdens
vloeden halen.
\end{solution}

\begin{solution}{exo:b1:clastic-sediments:10}
(a)~Met de middelpunten $-0.5, 0.5, \ldots, 4.5$: $m = (-1 + 4 + 45 + 100 + 52.5
+ 22.5)/100 = 2.23\,\varphi$; de gekwadrateerde afwijkingen, gewogen met de
percentages, tellen op tot \num{107.71}, dus $s = \sqrt{1.0771} = 1.04\,\varphi$:
slecht gesorteerd. (b)~Het cumulatieve percentage grover is 0, 2, 10, 40, 80, 95 en
100 bij $-1, 0, \ldots, 5\,\varphi$. Dan is $\varphi_5 = 0 + 3/8 = 0.375$,
$\varphi_{16} = 1 + 6/30 = 1.20$, $\varphi_{50} = 2 + 10/40 = 2.25$,
$\varphi_{84} = 3 + 4/15 = 3.27$, $\varphi_{95} = 4.0$. $M_z = (1.20 + 2.25 +
3.27)/3 = 2.24\,\varphi$ en $\sigma_I = (3.27 - 1.20)/4 + (4.0 - 0.375)/6.6 =
0.517 + 0.549 = 1.07\,\varphi$. (c)~De twee methoden komen tot op ongeveer
\num{0.01} overeen voor het gemiddelde en \num{0.03} voor de \omterm{def:b1:clastic-sediments:sorting}{sortering}: een fijn zand
($2.24\,\varphi$ is \qty{0.21}{mm}), slecht gesorteerd.
\end{solution}

\begin{solution}{exo:b1:clastic-sediments:11}
(a)~$U = \sqrt{0.02 \times 9.81 \times 50 \times 0.005/0.003} =
\qty{4.04}{m.s^{-1}}$, of 7.9 knopen. (b)~$(\Delta\rho/\rho)\,h = U^2
C_f/(gS) = 25.7^2 \times 0.003/(9.81 \times 0.005) = \qty{40}{m}$: bijvoorbeeld
een stroming van \qty{400}{m} dik met een overdichtheid van \qty{10}{\percent}, wat laat zien
hoe groot en dicht de stroming van 1929 moet zijn geweest. (c)~De snelheid schaalt als
$\sqrt S$: $4.04/\sqrt 5 = \qty{1.81}{m.s^{-1}}$. De langzamere stroming heeft
zwakkere turbulentie en laat haar grovere korrels vallen; zij wordt minder dicht, vertraagt
verder, en spreidt haar vracht over de vlakte uit als een \omterm{def:b1:clastic-sediments:turbidite}{gegradeerde laag}.
\end{solution}

\begin{solution}{exo:b1:clastic-sediments:12}
(a)~$0.7 \times \qty{5e10}{kg}/\qty{2000}{kg.m^{-3}} = \qty{1.75e7}{m^3}$ per
jaar. (b)~Uitgespreid over een front van \qty{30}{m} hoog en \qty{100}{km} breed, schuift
het front op met $\num{1.75e7}/(30 \times \num{1e5}) = \qty{5.8}{m}$ per jaar. (c)~$a = v\tan\alpha = 5.83 \times \tan\ang{0.5} = \qty{0.051}{m}$, ongeveer
\qty{5}{cm} per jaar, honderd keer minder dan de voortgang; de \qty{30}{m} opeenvolging
vormt zich in $30/0.051 = 590$ jaar.
\end{solution}

\begin{solution}{pb:b1:clastic-sediments:1}
\textbf{1.} Eenheid~1: $(5.0 + 6.6 + 8.4)/3 = 6.67\,\varphi$, \qty{0.0098}{mm}
(\qty{9.8}{\micro m}); eenheid~2: $3.47\,\varphi$, \qty{0.090}{mm}; eenheid~3:
$1.80\,\varphi$, \qty{0.29}{mm}; eenheid~4: $0.73\,\varphi$, \qty{0.60}{mm}.

\textbf{2.} Eenheid~1: $(8.4 - 5.0)/4 + (9.6 - 4.2)/6.6 = 0.850 + 0.818 =
1.67\,\varphi$; eenheid~2: $0.550 + 0.545 = 1.10$; eenheid~3: $0.275 + 0.273 = 0.55$;
eenheid~4: $0.725 + 0.712 = 1.44$.

\textbf{3.} Eenheid~1: fijne silt, slecht gesorteerd; eenheid~2: zeer fijn zand, slecht
gesorteerd; eenheid~3: middelgrof zand, matig goed gesorteerd; eenheid~4: grof zand,
slecht gesorteerd.

\textbf{4.} Eenheid~3. Aan de riviermonding bewerken de stroming en de golven het zand
steeds opnieuw en voeren zij de fijne delen weg, terwijl de modder van de prodelta uit
suspensie bezinkt als een mengsel van alle fijne groottes die haar bereiken.

\textbf{5.} De grootte $8\,\varphi$ ligt $1.4/1.8$ van de weg van
$\varphi_{50}$ naar $\varphi_{84}$, dus $50 + (1.4/1.8) \times 34 = 76$ procent van het
monster is grover en ongeveer \qty{24}{\percent} is van kleigrootte. Zulke fijne deeltjes
bezinken alleen uit bijna stil water: de plaats lag onder het bereik van golven en
stromingen.

\textbf{6.} Een geul voert zowel grind, dat in vloeden over haar bodem rolt, als zand; de
afzetting mengt het grindrestant van de diepste stroming met het zand van lagere
stromingen, zodat zij vele klassen bestrijkt.

\textbf{7.} Het gemiddelde groeit van \qty{0.0098}{mm} tot \qty{0.090}{mm} en
\qty{0.29}{mm}: de opeenvolging wordt naar boven grover.

\textbf{8.} Eenheid~1, gelamineerde modder: de prodelta; eenheid~2, afwisselend silt en
zand: het onderste deltafront; eenheid~3, kruisgelaagd zand: de monding-banken van het
bovenste deltafront.

\textbf{9.} Elke zone kwam te liggen op die welke eerder vóór haar lag: de delta bouwde
vooruit, zeewaarts, over de plaats heen.

\textbf{10.} Minstens \qty{21}{m}, de dikte van de eenheden 1 tot 3, omdat de opeenvolging
het water vulde van de bodem van de prodelta tot de monding-bank dicht bij het oppervlak;
meer, omdat de modders sindsdien zijn gecompacteerd
(\cref{ch:b1:diagenesis-basins}).

\textbf{11.} Een distributaire geul van de deltavlakte, die in de monding-bank sneed
terwijl de delta vooruit bouwde; haar kiezels zijn haar \omterm{def:b1:erosion-transport:load}{bodemvracht}.

\textbf{12.} Nee: een puntbank begint met een erosievlak en een restant en wordt naar boven
fijner, tot uiterwaardemodder. Het deltafront wordt van modder af naar boven grover.

\textbf{13.} $\qty{16}{m}/\num{4000}\text{ jaar} = \qty{4.0}{mm}$ per jaar.

\textbf{14.} $v = a/\tan\alpha = 0.0040/\tan\ang{0.5} = 0.0040/0.008727 =
\qty{0.458}{m}$ per jaar.

\textbf{15.} $0.458 \times 4000 = \qty{1830}{m}$, ongeveer \qty{1.8}{km}.

\textbf{16.} $\qty{21}{m}/\qty{4.0}{mm.yr^{-1}} = 5250$ jaar.

\textbf{17.} Het deltafront ontvangt het meeste zand nabij de monden, het prodeltagebied
veel minder en de deltavlakte weinig behalve bij vloeden; de actieve monden verschuiven ook
van plaats tot plaats, zodat een punt een tijd veel kan krijgen en daarna bijna niets. Een
enkele snelheid tussen twee gedateerde lagen is een gemiddelde over deze veranderingen, en
de modders hebben door compactie een deel van hun dikte verloren.

\textbf{18.} De som is 1544, dus het rekenkundig gemiddelde is \ang{128.7}, een
zuidoostelijke richting: zinloos, omdat de metingen rond het noorden verspreid liggen en
de waarden nabij 340 en nabij 010 alsof zij ver uit elkaar liggen worden gemiddeld.

\textbf{19.} $S = \sum\sin\alpha_i = 1.698$ en $C = \sum\cos\alpha_i =
11.248$.

\textbf{20.} $\bar R = \sqrt{1.698^2 + 11.248^2}/12 = 11.376/12 = 0.948$: de metingen zijn
strak rond één richting gegroepeerd.

\textbf{21.} De rivieren stroomden, en de delta bouwde, naar ongeveer het noorden
(\cref{prop:b1:clastic-sediments:azimuths}); het land lag in het zuiden.

\textbf{22.} Omdat de voorlagen zijn gemeten op lagen die niet meer horizontaal zijn: hun
\omterm{def:b1:geological-maps:attitude}{hellingrichtingen} zijn die van de gekantelde lagen plus de kanteling. Door elke voorlaag
om de \omterm{def:b1:geological-maps:attitude}{strijk} van de laagdeling te draaien over de helling van de laagdeling, wordt de
laagdeling, en daarmee de voorlagen, in hun oorspronkelijke ligging teruggebracht.

\textbf{23.} Getijstromingen stromen in en uit en bouwen sets die in twee tegengestelde
richtingen helen, zoals in visgraatkruislagen; het rozendiagram zou twee lobben hebben op
ongeveer \ang{180} van elkaar. De twaalf metingen vormen één lob, $\bar R$ nabij 1: een
rivierstroming die één kant op stroomt.

\textbf{24.} $\bar\alpha = \operatorname{atan2}(1.698, 11.248) = \ang{8.6}$: het gemiddelde
paleostroomazimut van eenheid~3 is \ang{9}, iets ten oosten van het noorden.
\end{solution}
